\section*{Lý do chọn đề
tài}

Trong quá trình chuyển đổi số lĩnh vực hành chính công, nhu cầu tra cứu
thủ tục của người dân ngày càng tăng. Người dân thường cần những thông
tin rất cụ thể như hồ sơ phải chuẩn bị, cơ quan tiếp nhận, thời hạn xử
lý, phí hoặc lệ phí. Tuy nhiên, các thông tin này thường được trình bày
trong những trang thủ tục dài, có nhiều thuật ngữ chuyên môn và yêu cầu
người dùng biết khá rõ tên thủ tục trước khi tìm kiếm.

Với tỉnh Lai Châu, đặc điểm địa bàn miền núi, dân cư phân tán và mức độ
tiếp cận công nghệ chưa đồng đều làm cho việc hỗ trợ tra cứu trực tuyến
trở nên cần thiết hơn. Một công cụ hỏi đáp bằng tiếng Việt tự nhiên có
thể giúp người dân kiểm tra trước những thông tin cơ bản, giảm số lần đi
lại và hạn chế tình trạng chuẩn bị thiếu giấy tờ khi đến cơ quan tiếp
nhận.

Đề tài lựa chọn kỹ thuật Retrieval-Augmented Generation (RAG) vì lĩnh
vực hành chính công yêu cầu câu trả lời phải bám theo dữ liệu nguồn. Mô
hình ngôn ngữ lớn (LLM) có thể diễn đạt tốt, nhưng nếu không có cơ chế
truy xuất và ràng buộc dữ liệu thì dễ phát sinh thông tin không có căn
cứ. RAG cho phép hệ thống tìm tài liệu liên quan trước, sau đó dùng mô
hình ngôn ngữ để diễn giải lại thông tin theo cách dễ đọc hơn.

Từ những lý do trên, đề tài ``Nghiên cứu và xây dựng chatbot hỗ trợ tra
cứu thủ tục hành chính (TTHC) cho tỉnh Lai Châu dựa trên kỹ thuật RAG''
được thực hiện nhằm xây dựng một hệ thống hỏi đáp trực tuyến, phục vụ
nhu cầu tra cứu thông tin thủ tục hành chính ở mức tư vấn ban đầu.

\section*{Mục đích nghiên
cứu}

Mục tiêu tổng quát của đề tài là xây dựng một ứng dụng chatbot web có
khả năng tiếp nhận câu hỏi bằng ngôn ngữ tự nhiên, truy xuất thông tin
từ kho dữ liệu TTHC và sinh câu trả lời ngắn gọn, rõ ràng, có kiểm soát
theo dữ liệu nguồn.

\begin{itemize}
\item
  Nghiên cứu tổng quan về chatbot, LLM, embedding, vector database và kỹ
  thuật RAG.
\item
  Khảo sát bài toán tra cứu TTHC, xác định nhóm thông tin người dân
  thường hỏi và yêu cầu quản trị dữ liệu.
\item
  Thu thập, chuẩn hóa và tổ chức dữ liệu TTHC tỉnh Lai Châu thành dạng
  phù hợp cho hệ thống RAG.
\item
  Xây dựng pipeline gồm các bước: làm sạch dữ liệu, chia chunk, tạo
  embedding, lưu vào ChromaDB, truy xuất, reranking và sinh câu trả lời.
\item
  Triển khai backend API, giao diện web chatbot, chức năng xác thực, lưu
  lịch sử hội thoại và trang quản trị dữ liệu.
\item
  Kiểm thử hệ thống theo các nhóm tiêu chí: chức năng, bảo mật, chất
  lượng truy xuất, độ đúng của câu trả lời và thời gian phản hồi.
\end{itemize}

\section*{Đối tượng và phạm vi nghiên
cứu}

Đối tượng nghiên cứu của đề tài gồm dữ liệu TTHC, quy trình người dân
tra cứu thông tin, kỹ thuật xử lý ngôn ngữ tự nhiên (NLP), mô hình RAG,
cơ sở dữ liệu vector và kiến trúc ứng dụng web client-server.

Phạm vi hệ thống tập trung vào chức năng hỗ trợ tra cứu thông tin TTHC
cho tỉnh Lai Châu thông qua giao diện chatbot web. Hệ thống xử lý câu
hỏi dạng văn bản tiếng Việt; chưa triển khai nhận dạng giọng nói, OCR
giấy tờ, thanh toán trực tuyến, ký số hoặc nộp hồ sơ trực tiếp lên hệ
thống nghiệp vụ của cơ quan nhà nước.

\textbf{Nguồn dữ liệu sử dụng trong đồ án được thu thập từ Cổng Dịch vụ
công tỉnh Lai Châu và đối chiếu với Cổng Dịch vụ công Quốc gia. Trong
quá trình xây dựng bộ dữ liệu thử nghiệm, một số mã thủ tục và mã định
danh được đối chiếu theo dữ liệu công khai liên quan đến TTHC. Dữ liệu
được chốt trong giai đoạn thực hiện đồ án tháng 5 năm 2026; khi triển
khai chính thức cần rà soát lại theo thời điểm vận hành thực tế vì TTHC
có thể thay đổi theo văn bản mới.}

Định dạng dữ liệu ban đầu là nội dung thủ tục hiển thị trên cổng dịch vụ
công dưới dạng trang thông tin/văn bản theo từng thủ tục. Sau khi thu
thập, dữ liệu được chuyển về dạng JSON để hệ thống dễ xử lý. Mỗi bản ghi
gồm mã thủ tục, tên thủ tục và phần nội dung chi tiết như lĩnh vực, đối
tượng thực hiện, cơ quan thực hiện, cách thức thực hiện, thời hạn, phí,
lệ phí, thành phần hồ sơ và căn cứ pháp lý.

Bộ dữ liệu thử nghiệm hiện có 1.950 TTHC, thuộc 176 lĩnh vực. Sau bước
kiểm tra, bộ dữ liệu không còn trùng mã thủ tục hoặc trùng tên thủ tục ở
mức bản ghi chính. Những trường rỗng, giá trị ``không quy định'', nội
dung thừa từ giao diện cổng và các đoạn nhiễu được chuẩn hóa trước khi
đưa vào pipeline RAG. Các thông tin dài được chia thành chunk theo ý
nghĩa nghiệp vụ để phục vụ tìm kiếm ngữ nghĩa.

\section*{Phương pháp nghiên
cứu}

Đề tài kết hợp giữa nghiên cứu lý thuyết và thực nghiệm phần mềm. Phần
lý thuyết dùng để xác định nền tảng công nghệ; phần thực nghiệm dùng để
kiểm tra các lựa chọn kỹ thuật trong bối cảnh dữ liệu TTHC tiếng Việt.

\begin{itemize}
\item
  Phương pháp khảo sát và phân tích nghiệp vụ: xác định quy trình tra
  cứu, tác nhân sử dụng, nhu cầu thông tin và phạm vi hệ thống.
\item
  Phương pháp xử lý dữ liệu: chuẩn hóa văn bản, loại bỏ nội dung nhiễu,
  kiểm tra trùng lặp, tách dữ liệu theo trường nghiệp vụ và chuyển thành
  JSON thống nhất.
\item
  Phương pháp thiết kế hệ thống: xây dựng Use Case, thiết kế kiến trúc
  backend/frontend, dữ liệu Firestore, metadata ChromaDB và pipeline
  RAG.
\item
  Phương pháp thực nghiệm: triển khai mã nguồn, chạy thử với dữ liệu
  thật, quan sát log truy xuất và điều chỉnh chunking, prompt hoặc
  reranking khi có lỗi.
\item
  Phương pháp kiểm thử: xây dựng bộ câu hỏi thử nghiệm theo nhóm hồ sơ,
  phí/lệ phí, thời hạn, cơ quan thực hiện, câu hỏi tiếp nối và câu hỏi
  ngoài phạm vi.
\end{itemize}

\section*{Cấu trúc báo
cáo}

Báo cáo được trình bày theo bốn chương, gồm:

\begin{itemize}
\item
  Chương 1. CƠ SỞ LÝ THUYẾT
\item
  Chương 2. KHẢO SÁT VÀ PHÂN TÍCH HỆ THỐNG
\item
  Chương 3. THIẾT KẾ HỆ THỐNG
\item
  Chương 4. TRIỂN KHAI
\end{itemize}
